\section*{Chapitre \ref{ch:b3:gfd-foundations} --- Dynamique des fluides géophysiques}

\begin{solution}{exo:b3:gfd-foundations:1}
(a)~$f = 2\Omega\sin\varphi$~: \num{7.292e-5}, \num{1.031e-4} et
\qty{1.263e-4}{s^{-1}}~; $2\pi/f$~: 23.93, 16.92 et \qty{13.82}{h}.
(b)~$U/f$~: 2.74, 1.94 et \qty{1.58}{km}. (c)~À \ang{30}, $\sin\varphi =
1/2$ et $2\pi/f = 2\pi/\Omega$~: la période d'inertie est le jour sidéral,
\qty{23.93}{h}.
\end{solution}

\begin{solution}{exo:b3:gfd-foundations:2}
Avec $f = \qty{1.031e-4}{s^{-1}}$~: tourbillon de poussière \num{9.7e3}, brise de mer 2.4,
dépression 0.15, gyre \num{2.4e-4}. La dépression et le gyre sont proches de
l'équilibre géostrophique~; le tourbillon de poussière tourne dans un sens ou dans l'autre, sans tenir compte de la
rotation terrestre, et la brise de mer ne la ressent qu'au bout d'une journée.
\end{solution}

\begin{solution}{exo:b3:gfd-foundations:3}
$\partial u_g/\partial z = (g/fT)\,|\partial T/\partial y| = 9.81\times
\num{5e-6}/(\num{1.031e-4}\times250) = \qty{1.90e-3}{s^{-1}}$~; sur
\qty{10}{km}, le vent d'ouest augmente de \qty{19}{m/s}.
\end{solution}

\begin{solution}{exo:b3:gfd-foundations:4}
$E = \num{50e-8}$~: $N = \qty{2.21e-3}{s^{-1}}$, une période de \qty{47}{min},
\num{1.27} cycles par heure. $E = \num{1000e-8}$~: $N = \qty{9.90e-3}{s^{-1}}$,
\qty{10.6}{min}, 5.7 cycles par heure. $E = \num{e-8}$~: $N =
\qty{3.13e-4}{s^{-1}}$, \qty{5.6}{h}, 0.18 cycles par heure.
\end{solution}

\begin{solution}{exo:b3:gfd-foundations:5}
(a)~$W \sim UH/L = \qty{0.1}{m/s}$ et $\mathrm Dw/\mathrm Dt \sim UW/L =
\qty{e-6}{m/s^2}$, \num{1.0e-7} de $g$. Le terme de pression, $fUL/H =
\qty{0.103}{m/s^2}$, est supérieur d'un facteur \num{1e5}~: le rapport, \num{9.7e-6}, est
$\mathrm{Ro}(H/L)^2$ avec $\mathrm{Ro} = 0.097$. (b)~$\tilde fu =
2\Omega\cos\ang{45}\times10 = \qty{1.03e-3}{m/s^2}$, \num{1.05e-4} de $g$.
\end{solution}

\begin{solution}{exo:b3:gfd-foundations:6}
$fr = \qty{30}{m/s}$. (a)~$\mathrm{Ro}_g = 0.667$, $V = 40/(1 + \sqrt{3.667})
= \qty{13.7}{m/s}$, un tiers plus lent que le vent géostrophique. (b)~$\mathrm{Ro}_g =
0.167$, $V = 10/(1 + \sqrt{0.333}) = \qty{6.34}{m/s}$. (c)~$V_g \le fr/4 =
\qty{7.5}{m/s}$, et alors $V = fr/2 = \qty{15}{m/s}$, soit deux fois
la vitesse géostrophique.
\end{solution}

\begin{solution}{exo:b3:gfd-foundations:7}
(a)~$\Omega = 2\pi/6 = \qty{1.047}{rad/s}$, $f = \qty{2.094}{s^{-1}}$.
$\mathrm{Ro} = 0.01/(2.094\times0.05) = 0.095$ et $\mathrm{Ek} =
\num{e-6}/(2.094\times0.04) = \num{1.19e-5}$~: les deux sont petits, donc l'écoulement obéit
au \omterm{thm:b3:gfd-foundations:taylor-proudman}{théorème de Taylor--Proudman}, et une \omterm{thm:b3:gfd-foundations:taylor-proudman}{colonne de Taylor} se forme. (b)~$0.2/\sqrt{\num{e-6}
\times1.047} = \qty{195}{s}$, \qty{3.3}{min}, soit 33 tours du plateau.
\end{solution}

\begin{solution}{exo:b3:gfd-foundations:8}
$\partial\bar T/\partial y = -f\,\Delta u/(R\ln(850/250))$, avec
$\ln(850/250) = 1.224$. À \ang{30} N, $-\num{7.29e-5}\times36.12/(286.9
\times1.224) = \qty{-7.50}{K}$ par \qty{1000}{km}~; à \ang{45} N,
$-\num{1.031e-4}\times17.18/351.1 = \qty{-5.05}{K}$ par \qty{1000}{km}. Le
cisaillement est proportionnel au gradient divisé par $f$~: le gradient est 1.5
fois plus grand à \ang{30} et $f$ 1.4 fois plus petit, donc le cisaillement est 2.1 fois
plus grand.
\end{solution}

\begin{solution}{exo:b3:gfd-foundations:9}
(a)~$L_R = 0.0113\times11\,000/\num{1.031e-4} = \qty{1206}{km}$;
$0.00357\times1000/\num{9.47e-5} = \qty{37.7}{km}$, 32 fois plus petit.
(b)~$\beta = 2\Omega\cos\varphi/a$~: \num{1.62e-11} et
\qty{1.74e-11}{m^{-1}.s^{-1}}. $\beta L_R/f = 0.19$ pour l'atmosphère, 0.007
pour l'océan.
\end{solution}

\begin{solution}{exo:b3:gfd-foundations:10}
(a)~$\mathrm dw/\mathrm dt = -(r + \mathrm if)w$, donc $w = U\mathrm
e^{-rt}\mathrm e^{-\mathrm ift}$~: la vitesse tourne au taux $f$ tandis que
son module décroît comme $\mathrm e^{-rt}$. La position $x + \mathrm iy =
U(1 - \mathrm e^{-(r + \mathrm if)t})/(r + \mathrm if)$ s'enroule vers le
point $U/(r + \mathrm if)$. (b)~$x + \mathrm iy = U(1 - \mathrm e^{-(r +
\mathrm if)t})/(r + \mathrm if) + u_0t$~: la spirale s'étire en boucles.
(c)~$0.2\,\mathrm e^{-1} = \qty{0.074}{m/s}$~; le centre a dérivé de $u_0t =
0.05\times4\times86\,400 = \qty{17.3}{km}$ vers l'est.
\end{solution}

\begin{solution}{exo:b3:gfd-foundations:11}
$f = \qty{9.37e-5}{s^{-1}}$; $\partial u/\partial z = (g/f\rho_0)\,\partial
\rho/\partial y = 9.81\times10^{-5}/(\num{9.37e-5}\times1025) =
\qty{1.02e-3}{s^{-1}}$. Avec $u = 0$ à \qty{1000}{m}, le courant de surface est
\qty{1.02}{m/s}, vers l'est~: dans l'hémisphère Nord, le courant s'écoule avec
l'eau légère à sa droite.
\end{solution}

\begin{solution}{exo:b3:gfd-foundations:12}
(a)~$N^2 = (g/\rho)\,\Delta\sigma_\theta/\Delta z$. 0--\qty{100}{m}~:
$9.81\times0.668/(1026.4\times100)$, $N = \qty{7.99e-3}{s^{-1}}$, 4.6
cycles par heure~; 100--\qty{500}{m}~: \qty{3.02e-3}{s^{-1}}, 1.7 cycles par
heure~; 500--\qty{1000}{m}~: \qty{3.24e-3}{s^{-1}}, 1.9 cycles par heure.
(b)~La tranche supérieure \qty{100}{m}~: elle contient l'eau chaude de surface au-dessus de la
pycnocline saisonnier. Entre 100 et \qty{500}{m}, l'eau est plus
uniforme qu'en dessous. (c)~$E = N^2/g$ vaut \num{650e-8}, \num{93e-8} et
\qty{107e-8}{m^{-1}}, toutes dans la plage de valeurs de l'océan supérieur.
\end{solution}

\begin{solution}{pb:b3:gfd-foundations:1}
\textbf{1.} $f = 2\times\num{7.292115e-5}\sin\ang{48} =
\qty{1.084e-4}{s^{-1}}$; $2\pi/f = \qty{16.1}{h}$.

\textbf{2.} Baignoire \num{3.1e3}~; cuve ($f = 4\pi/10 = \qty{1.257}{s^{-1}}$)
0.020~; brise 2.3~; dépression 0.14~; gyre \num{2.3e-4}.

\textbf{3.} La baignoire ne ressent rien de la rotation terrestre~; la brise
ne la ressent qu'après quelques heures. La cuve, la dépression et le gyre ont
de petits \omterm{def:b3:gfd-foundations:numbers}{nombres de Rossby} et sont proches de l'équilibre géostrophique.

\textbf{4.} Cuve $\num{e-6}/(1.257\times0.04) = \num{2.0e-5}$~; dépression
$10/(\num{1.084e-4}\times10^8) = \num{9.2e-4}$~; gyre
$\num{e-3}/(\num{1.084e-4}\times10^6) = \num{9.2e-6}$.

\textbf{5.} $\delta/H = \sqrt{2\nu/(fH^2)} = \sqrt{2\,\mathrm{Ek}}$. Cuve
$\sqrt{2\times\num{e-6}/1.257} = \qty{1.3}{mm}$~; air $\sqrt{20/\num{1.084e-4}}
= \qty{430}{m}$.

\textbf{6.} Le frottement agit dans des couches de quelques millièmes de la profondeur dans la
cuve et dans l'océan, de quelques centièmes dans l'atmosphère~; au-dessus d'elles, l'écoulement
est presque sans frottement.

\textbf{7.} $\rho = p/(RT) = 101\,290/(286.9\times288.1) =
\qty{1.225}{kg/m^3}$.

\textbf{8.} $u_g = (1/\rho f)\,|\partial p/\partial y| = (400/300\,000)/(1.225
\times\num{1.084e-4}) = \qty{10.0}{m/s}$, venant de l'ouest, la basse pression
étant au nord, à gauche.

\textbf{9.} $\mathrm{Ro}_g = 10.04/(\num{1.084e-4}\times600\,000) = 0.154$;
$V = 2\times10.04/(1 + \sqrt{1.618}) = \qty{8.84}{m/s}$, \qty{12}{\percent}
plus lent.

\textbf{10.} $V = 2\times10.04/(1 + \sqrt{0.382}) = \qty{12.4}{m/s}$. La
plus grande $V_g$ est $fr/4 = \qty{16.3}{m/s}$, un gradient $\rho fV_g =
\qty{2.16e-3}{Pa/m}$, soit \qty{6.5}{hPa} par \qty{300}{km}.

\textbf{11.} $g/c_p = \qty{9.77}{K/km}$. $N^2 = (9.81/255.7)\times(9.77 -
6.49)\times10^{-3} = \qty{1.26e-4}{s^{-2}}$, $N = \qty{0.0112}{s^{-1}}$, soit
une période de \qty{9.3}{min}.

\textbf{12.} Avec $\Gamma = 0$, $N^2 = (g/T)(g/c_p) = g^2/(c_pT) =
\qty{4.43e-4}{s^{-2}}$, $N = \qty{0.0210}{s^{-1}}$, soit 1.9 fois la
valeur troposphérique.

\textbf{13.} $L_R = 0.0112\times10^4/\num{1.084e-4} = \qty{1035}{km}$, la
taille de la dépression.

\textbf{14.} $g' = \qty{0.0147}{m/s^2}$; $\sqrt{g'h} = \qty{0.54}{m/s}$;
$L_R = 0.542/\num{1.084e-4} = \qty{5.0}{km}$.

\textbf{15.} Oui~: le lac, large de \qty{12}{km}, fait plus du double
de son rayon de Rossby, si bien que ses ondes internes et ses courants sont déviés par la
force de Coriolis.

\textbf{16.} $(249.4 - 253.0)/600 = \qty{-6.0}{K}$ par \qty{1000}{km}.

\textbf{17.} $\ln(1013/265) = 1.3409$. Sud~:
$286.9\times253.0\times1.3409/9.81 = \qty{9922}{m}$~; nord \qty{9781}{m}.
La surface à \qty{265}{hPa} est \qty{141}{m} plus basse au nord, une pente de
\num{2.35e-4} par rapport à la surface inférieure.

\textbf{18.} $\Delta u = (g/f)\times\text{pente} = 9.81\times\num{2.35e-4}/
\num{1.084e-4} = \qty{21.3}{m/s}$.

\textbf{19.} $(R/f)\ln(p_b/p_t)\,|\partial\bar T/\partial y| =
(286.9/\num{1.084e-4})\times1.3409\times\num{6.0e-6} = \qty{21.3}{m/s}$, la
même~: les deux calculs ne font qu'un.

\textbf{20.} \omterm{def:b3:gfd-foundations:baroclinic}{Barocline}~: la température, donc la masse volumique, varie le long
des surfaces de pression. \omterm{def:b3:gfd-foundations:baroclinic}{Barotrope}, le vent ne changerait pas avec l'altitude~:
\qty{10.0}{m/s}.

\textbf{21.} $(g/f\bar T)\,|\partial T/\partial y|\,\Delta z = 9.81\times
\num{6.0e-6}\times10^4/(\num{1.084e-4}\times255.7) = \qty{21.2}{m/s}$,
à \qty{0.3}{\percent} près.

\textbf{22.} $10.04 + (286.9/\num{1.084e-4})\times\num{6.0e-6}\times
\ln(1013/541) = 10.04 + 9.96 = \qty{20.0}{m/s}$.

\textbf{23.} $\partial u/\partial z = 21.3/10^4 = \qty{2.13e-3}{s^{-1}}$;
$\mathrm{Ri} = \num{1.26e-4}/\num{4.54e-6} = 28$, bien au-dessus de 0.25~: le cisaillement
ne rend pas turbulente la troposphère libre.

\textbf{24.} Le vent d'Orrin, \qty{31}{m/s}, est plus fort que la moyenne zonale
à \ang{45} N et plus faible qu'à \ang{30} N. Le courant-jet du front polaire oscille
vers le nord et vers le sud~; à tout instant, son gradient et son vent sont concentrés
sur une bande, et la moyenne sur un cercle de latitude l'étale.

\textbf{25.} $10.0 + 21.3$~: le vent géostrophique à \qty{10}{km} au-dessus de la
steppe est \textbf{\qty[text-series-to-math=true]{31}{m/s}}, venant de l'ouest
(\qty{31.3}{m/s} par le calcul).
\end{solution}
